\chapter*{Foreword}
\addcontentsline{toc}{chapter}{Foreword}
\markboth{Foreword}{}

As I write this, I am in awe of the tools available to anyone with an idea, the freedom to pursue it, and the commitment to keep learning. AI is accelerating what people are discovering and building. Powering that progress raises a question that matters to every household: how can new large electricity customers, like Bitcoin miners and AI data centers, help make power more reliable and affordable for the people already on the grid?

Bitcoin miners have spent years learning part of the answer. Inspired by the oil and gas wildcatters who came before us, we've searched the far reaches of the earth for latent, cheap, and abundant electricity. In Texas, we encountered an opportunity made possible by pioneering people who had been thinking about these questions long before Bitcoin existed. ERCOT's market designers recognized the value of customers responding to prices that reflect when and where electricity is abundant or scarce. Their work created incentives to buy available power and reduce consumption as grid conditions tighten. Bitcoin miners brought a new kind of large, highly elastic load customer to that market framework.

What we came to understand was how pursuing affordable power could also benefit others. With the right market rules, our power purchases help support the economics of generation everyone depends on, while turning down leaves more electricity available when families, hospitals, and other essential services need it most. We began to see the potential for a virtuous cycle of mutual benefit, one still in its infancy.

Putting that opportunity to work required learning a new language and A LOT of acronyms. When we began engaging with ERCOT stakeholders in 2017, I struggled to understand grid protocols, explain our capabilities, and ask the right questions. We could see the potential of load as a resource, but we still had to learn how to make our flexibility something operators could count on through protocol revisions. Vit Ficl has written the book I wish I'd had.

Winter Storm Uri gave that work urgency. The hardship Texans endured in February 2021 reinforced how much families and businesses depend on reliable electricity. I remain grateful to the operators, lineworkers, and crews who serve our communities through difficult conditions, and committed to learning how we can better prepare for the next emergency. As a grid nerd who enjoys hard problems, I join those conversations eagerly, listening and learning. Understanding what went wrong is essential. So is asking, ``What worked, and how do we expand what worked?''

Bitcoin miners were among the largest customers who voluntarily reduced electricity consumption during Uri, helping ease demand on a strained system. Their response illustrated an opportunity worth developing: customers that can reliably turn down give grid operators another tool when power is scarce. Alongside dependable generation, battery storage, and resilient infrastructure, that flexibility can help strengthen reliability with the right regulatory framework. It is a lesson grids around the world can build on, and the potential contribution begins long before an emergency.

You have probably read that a Bitcoin mine consumes ``enough electricity to power thousands of homes.'' That comparison paints a frightening picture of computers taking power away from families. It leaves out an important part of the story: power plants need enough revenue throughout the year to stay in business. When those homes need less electricity, miners can buy available power, helping underused generators run more often and cover their costs. When a winter storm arrives and families need more electricity, miners can turn down. The same customer accused of taking electricity from thousands of homes is actually helping pay to keep generation available, and then stepping aside so it can serve those same homes when they need it most.

That is the opportunity for mutual benefit: affordable power for miners, revenue for generators, and economic support for the dependable supply families, hospitals, and critical infrastructure need. Making that benefit dependable takes sound engineering, appropriate incentives, protocol revisions where needed, and performance grid operators can verify. This is where Vit's work becomes valuable. He connects what a Bitcoin mining server can do with what an electricity system needs, explains the economics, and shows his assumptions so readers can test the reasoning and improve on what they learn.

The lessons reach beyond Bitcoin mining. AI companies are working through questions miners know well: where to find power, how quickly they can access it, and how to support the infrastructure their growth requires without shifting the bill to other customers. Their operating needs differ, but where flexibility is practical and market rules recognize it, it can offer a path to faster access to power. Bitcoin mining offers experience worth learning from.

Bitcoin miners came looking for cheap, abundant electricity. By bringing computing demand to a network of underutilized generation, we discovered how pursuing our commercial interests could also make underused resources more productive and help sustain the supply others depend on. Others will build on what we have learned in ways we have yet to imagine. That is why this book matters. Knowledge that took years to gather through lived experience is now available to anyone willing to learn and put it to work. Somewhere in these pages, a reader may find the missing piece that turns an idea into something useful.

Mankind's quest to understand the universe has no finish line. Powering that pursuit gives us an opportunity to build businesses that make electricity more reliable, affordable, and abundant for everyone, so families can live better, communities can prosper, and the next generation can accomplish more.

I'm confident we're at the beginning of an age of energy abundance, and there's room for more pioneers to help shape the world we want. Take what Vit has learned and put your own ideas to work. There's a whole lot left to discover, and we're just getting warmed up!

\vspace{1em}
\noindent Gideon O. Powell \\
CEO, Cholla Inc.
